
\noindent\textit{Proof.}
For every \(i,j\geq0\), put
\[
  W_{i,j}=U_i\widehat\otimes V_j,
  \qquad
  \overline W_{i,j}
  =
  (U_i/U_{i+1})\widehat\otimes_k(V_j/V_{j+1}).
\]
We prove by induction on \(n\) that the natural map
\[
  \pi_n:W_n\longrightarrow
  \bigoplus_{i+j=n}\overline W_{i,j}
\]
is surjective and that the inclusion
\[
  W_{n+1}\subset\ker(\pi_n)
\]
is an equality. For \(n=0\), the projection
\[
  \pi_0:
  U_0\widehat\otimes_kV_0
  \longrightarrow
  (U_0/U_1)\widehat\otimes_k(V_0/V_1)
\]
is surjective. Since \(U_0\), \(U_0/U_1\), and hence \(U_1\),
are topologically flat over \(k\), we have
\[
  \ker(\pi_0)
  =
  U_0\widehat\otimes_kV_1
  +
  U_1\widehat\otimes_kV_0,
\]
and moreover
\[
  \bigl(U_0\widehat\otimes_kV_1\bigr)
  \cap
  \bigl(U_1\widehat\otimes_kV_0\bigr)
  =
  U_1\widehat\otimes_kV_1.
\]

Suppose now that \(n>0\) and that the result has been proved for
\(n-1\). Put
\[
  M_0=U_0\widehat\otimes_kV_n,
  \qquad
  S_0=\sum_{i=1}^{n}U_i\widehat\otimes_kV_{n-i}.
\]
We have \(S_0\subset U_1\widehat\otimes_kV_0\). Applying the
preceding argument to \(V_0\supset V_n\) in place of
\(V_0\supset V_1\), we obtain
\[
  M_0\cap S_0
  \subset
  \bigl(U_0\widehat\otimes_kV_n\bigr)
  \cap
  \bigl(U_1\widehat\otimes_kV\bigr)
  =
  U_1\widehat\otimes_kV_n.
\]
It follows that
\[
  M_0\cap S_0=U_1\widehat\otimes_kV_n.
\]
Since \(W_n=M_0+S_0\), we obtain a commutative diagram with exact
rows, where
\[
  U'_i=U_{i+1},
  \qquad
  W'_{i,n-1-i}=W_{i+1,n-1-i}
  \quad(0\leq i\leq n-1):
\]
\begin{center}
  \phantomsection
  \label{sga3:VIIB:5.1.2:filtration-exact-rows-diagram}
  \begin{tikzcd}[row sep=3em,column sep=1.8em]
0 \arrow[r] & S_0 \arrow[r] \arrow[d,"\pi'_{n-1}"']
  & W_n \arrow[r] \arrow[d,"\pi_n"']
  & (U_0/U_1)\widehat\otimes_k V_n
      \arrow[r] \arrow[d,"p"']
  & 0 \\
0 \arrow[r]
  & \displaystyle\bigoplus_{i=0}^{n-1}\overline W'_{i,n-1-i}
      \arrow[r]
  & \displaystyle\bigoplus_{i=0}^{n}\overline W_{i,n-i}
      \arrow[r]
  & \overline W_{0,n} \arrow[r]
  & 0
\end{tikzcd}

\end{center}

\phantomsection
\label{sga3:VIIB:5.1.2:diagram-source-notation-note}
\noindent\textit{Translator's note.}\enspace
The French diagram's lower middle term is printed
\(\bigoplus_{i=0}^{n}\overline W_{i,n-1}\); the natural map
\(\pi_n\) defined immediately above instead has target
\(\bigoplus_{i+j=n}\overline W_{i,j}\), which would give
\(\bigoplus_{i=0}^{n}\overline W_{i,n-i}\). The complementary index is used in the diagram above.

Then \(p\) is surjective, with kernel
\[
  (U_0/U_1)\widehat\otimes_kV_{n+1}.
\]
Moreover, by the induction hypothesis applied to the sequence
\((U'_i)\), the map \(\pi'_{n-1}\) is surjective, with kernel
\[
  W'_n
  =
  \sum_{i=1}^{n}W_{i,n+1-i}.
\]
It follows that \(\pi_n\) is surjective and that
\[
  W_{n+1}\subset\ker(\pi_n)
\]
is an equality. This proves the lemma.

Let us return to an object \((A,J)\) of \(\mathcal C\). By
\SGARefLink{sga3:VIIB:0.2.G}{0.2.G}, the hypothesis that every
\[
  \overline{J^n}/\overline{J^{n+1}}
\]
is topologically flat over \(k\) is equivalent to saying that
\(\operatorname{Gr}_J(A)\) is topologically flat over \(k\).

\medskip
\noindent\textbf{Corollary.}
\phantomsection
\label{sga3:VIIB:5.1.2:corollary}
{\itshape Let \(\mathcal P\) be the full subcategory of
\(\mathcal C\) consisting of the objects \((A,J)\) such that
\(A\) and \(\operatorname{Gr}_J(A)\) are topologically flat over
\(k\). Then the functor
\[
  \mathcal P\longrightarrow\mathbf{Alpg}/k,
  \qquad
  (A,J)\longmapsto\operatorname{Gr}_J(A),
\]
commutes with finite coproducts and hence carries a cogroup of
\(\mathcal P\) to a cogroup of \(\mathbf{Alpg}/k\).

In particular, if \(k\) is a field, then for every cogroup
\((A,J)\) of \(\mathcal C\), the algebra
\(\operatorname{Gr}_J(A)\) is a cogroup of
\(\mathbf{Alpg}/k\).}
